\section{Conclusion}
\label{sec:conclusion}

We present SiliQun, an open-source quantum simulation platform that unifies silicon spin qubit technologies—SiMOS quantum dots, phosphorus-donor systems, and gate-all-around nanowires—within a physics-grounded Lindblad density-matrix framework and a Gymnasium-compatible reinforcement learning interface. Developed through the Physics-Grounded RL Simulation (PGIRS) methodology, SiliQun systematically integrates noise-aware design, hardware-profile abstraction, and empirical validation into a single reusable stack. It is, to our knowledge, the first open-source platform to satisfy all five jointly essential properties for silicon spin qubit RL research—Gymnasium compliance, Lindblad physics, hardware-calibrated parameters, gate-level synthesis, and open availability—validated through four independent replication studies across three hardware architectures and four RL algorithms.

The validation results demonstrate SiliQun's capabilities across several dimensions. First, the Lindblad physics engine achieves machine-precision accuracy in fundamental quantum benchmarks, including Rabi oscillations and $T_1$ decay processes, with maximum deviations of just $2.78 \times 10^{-16}$---effectively at the limit of float64 numerical precision. Second, our replication studies establish cross-technology consistency: we successfully reproduce Daraeizadeh~et~al.'s DQL/SARSA CZ gate synthesis on SiMOS platforms ($\bar{F} = 0.9996 \pm 0.0001$, within 0.04\% of the reference implementation) and He~et~al.'s universal state-preparation protocol on donor systems ($\bar{F} = 0.9748 \pm 0.0098$). Notably, when applying similar DQN protocols to gate-all-around nanowires, we obtain $\bar{F} = 0.9877 \pm 0.0021$, which we believe constitutes the first machine-learning benchmark for this emerging qubit platform.

Beyond these baseline validations, our reinforcement learning experiments reveal two complementary aspects of the system's physical realism. A lightweight PPO agent achieves near-perfect fidelity ($F = 0.9996$) on 2-qubit Bell state synthesis, confirming both the engine's learnable reward signals and correct Gymnasium implementation. However, the same agent manages only $F = 0.3998$ on the 4-qubit GHZ task, demonstrating physically realistic noise scaling and genuine algorithmic challenges absent specialised enhancements. These findings collectively suggest that SiliQun successfully balances pedagogical accessibility with research-grade physical accuracy.

Perhaps most revealing are our scalability studies, which expose fundamental limitations in current standard approaches. While SAC baselines reliably solve all 2-qubit target states ($F > 0.99$), performance sharply declines for 3-qubit GHZ, W, and cluster states ($\bar{F} = 0.69 \pm 0.06$). Interestingly, Dicke-$k_3$ states prove considerably more tractable ($\bar{F} = 0.9782 \pm 0.0137$ at $n=3$), hinting at state-dependent noise susceptibility driven by the lower entanglement entropy of this family. This scalability ceiling---which we analyse in depth through noise-conditional hyperparameter studies---directly motivated the development of a companion adaptive quantum control framework (under separate submission), where branched critic distillation and evolutionary reinforcement learning exploration overcome these limitations.

Architecturally, SiliQun makes two lasting contributions. Its hardware-neutral simulation substrate enables unprecedented cross-technology comparisons through a standardised Gymnasium interface, overcoming the tool fragmentation that has historically impeded silicon spin qubit research. The PGIRS methodology further provides transferable engineering principles for extending the platform to non-silicon qubit technologies without modifying the core Lindblad solver---a flexibility we have already begun exploiting in ongoing transmon research.

Three directions emerge as priorities for future work. Comparative studies between reinforcement learning and classical optimal control methods (GRAPE, CRAB) could clarify when machine learning offers practical advantages in gate synthesis. Physical validation through randomised benchmarking would help bridge the current simulation-to-reality gap. Finally, extending scalability studies beyond four qubits may reveal whether the observed noise-conditional hyperparameter relationships generalise to larger registers---a question we are actively investigating through PGIRS-guided design iterations.

SiliQun establishes a foundation for scalable silicon quantum computing control by unifying device-accurate simulation with reinforcement learning—enabling direct policy transfer from simulation to emerging qubit platforms such as GAA nanowires, where experimental benchmarks remain scarce. As fabrication advances bring GAA and donor-electron devices closer to operational use, the ability to pre-train noise-aware control policies in a physics-grounded simulator will become increasingly valuable. SiliQun is available under the MIT licence at \url{https://github.com/r3d-phd/siliqun}, with full reproducibility ensured through provided scripts and archived HPC job files. We invite community contributions to expand both its physical models and its pedagogical applications.
